
\section{Additional results for   \texorpdfstring{\Cref{chap:imitation_undelayed}}{}}


\subsection{Bounds Involving the Wasserstein Distance}


\begin{prop}
\label{pp:expected_wass_dist}
    Consider two real random variables $X,Y$ with respective distributions $\eta,\nu$. Then, one has:
    \begin{align*}
        \left\vert \expectedvalue[X]-\expectedvalue[Y]  \right\vert\leq \wassersteindistance(\eta\Vert \nu).
    \end{align*}
\end{prop}
\begin{proof}
    One has,
    \begin{align}
        \expectedvalue[X]-\expectedvalue[Y] 
        &= \int_{\realnumbers} x(\eta(x)-\nu(x))~dx
        \nonumber\\
        &\leq\sup_{\left\Vert f\right\Vert_L\leq 1}\left\vert\int_{\realnumbers} f(x)(\eta(x)-\nu(x))~dx\right\vert
        \label{eq:x_lip}\\
        &\leq \wassersteindistance(\eta\Vert \nu),\label{eq:recognise_wass}
    \end{align}
    where \Cref{eq:x_lip} follows since the identity is 1-LC and \Cref{eq:recognise_wass} is obtained by recognising the Wasserstein distance.
    The same reasoning can be applied to $\expectedvalue[Y]-\expectedvalue[X]$ and the symmetry of the Wasserstein distance concludes.
\end{proof}

The next result asserts that if one applies a $L$-LC function to two random variables, one gets two random variables with distribution whose Wasserstein distance is bounded by the original Wasserstein distance multiplied by a factor $L$.

\begin{prop}
\label{pp:wass_dist_lip_g}
    Consider two probability measures $\eta$ and $\nu$ over the metric space $\Omega$ and an $L_f$-LC function $f:\Omega\rightarrow \realnumbers$.
    For some random variable $X$ distributed according to $\eta$, we note $f_\eta$ the distribution of $f(X)$. We define similarly $f_\nu$.
    One then has,
    \begin{align*}
        \wassersteindistance(f_{\eta}\Vert f_{\nu})\leq L_f \wassersteindistance(\eta \Vert \nu).
    \end{align*}
\end{prop}
\begin{proof}
    We show the results directly:
    \begin{align}
        \wassersteindistance(f_{\eta}\Vert f_{\nu}) 
        &= \sup_{\left\Vert g\right\Vert_L\leq 1}\left\vert\int_{\mathbb R} g(x)(f_{\eta}(x)-f_{\nu}(x)) ~dx\right\vert
        \label{eq:def_wass_f}
        \\
        &=\sup_{\left\Vert g\right\Vert_L\leq 1}
        \left\vert\int_{\mathbb R} g(x)f_{\eta}(x)~dx-\int_{\mathbb R} g(x)f_{\nu}(x) dx\right\vert,\nonumber
        \\
        &= \sup_{\left\Vert g\right\Vert_L\leq 1}
        \left\vert\int_{\mathbb R} g(f(x))\eta(x) ~dx-\int_{\mathbb R} g(f(x)) \nu(x) dx\right\vert\label{eq:def_eta_nu}
        \\
        &= \sup_{\left\Vert g\right\Vert_L\leq 1}
        \left\vert \int_{\mathbb R} g(f(x))(\eta(x) - \nu(x)) ~dx\right\vert,\nonumber
    \end{align}  
    where \Cref{eq:def_wass_f} holds by definition of the Wasserstein distance and \Cref{eq:def_eta_nu} holds by the definition of $f_\eta$ and $f_\nu$.
    The conclusion follows from observing that $g(f(x))$ is $L_f$-LC by composition of a 1-LC and $L_f$-LC functions.
\end{proof}

\begin{prop}
\label{pp:expected_q_bound}
    Let $\markovdecisionprocess$ be an \gls{mdp} and $\pi$ a policy whose Q-function is $L_{Q}$-LC in the second argument.
    Then, for any two probability distributions $\eta,\nu$ over $\actionspace$, 
    \begin{align*}
        \left\vert \expectedvalue_{\substack{X\sim \eta\\Y\sim \nu}}[Q^{\pi}(s,X)-Q^{\pi}(s,Y)] \right\vert 
        \leq L_{Q} \wassersteindistance(\eta(\cdot)\Vert\nu(\cdot)).
    \end{align*}
\end{prop}
\begin{proof}
    For some $s\in\statespace$, let $f_{\eta}(X)$ and $f_{\nu}(X)$ be the distributions of $Q^{\pi}(s,X)$ and $Q^{\pi}(s,Y)$ respectively.
    First, by applying \Cref{pp:expected_wass_dist}, one gets,
    \begin{align*}
        \left\vert \expectedvalue_{\substack{X\sim \eta\\Y\sim \nu}}[Q^{\pi}(s,X)-Q^{\pi}(s,Y)] \right\vert 
        \leq  \wassersteindistance(g_{\eta}\Vert g_{\nu}).
    \end{align*}
    Then, applying \Cref{pp:wass_dist_lip_g}, one gets,
    \begin{align*}
         \wassersteindistance(f_{\eta}\Vert f_{\nu})
         \leq L_{Q} \wassersteindistance(\eta\Vert \nu).
    \end{align*}
\end{proof}


\begin{prop}
\label{pp:lip_dmdp}
    Let $\markovdecisionprocess$ be a $(L_P,L_r)$-LC \gls{mdp}.
    Let $\delayedmarkovdecisionprocess$ be a constantly $\delay$-delayed \gls{dmdp} built upon $\markovdecisionprocess$ and with transition function $\augmentedtransitionfunction$ and reward function $\augmentedexpectedrewardfunction$.
    Then, $\delayedmarkovdecisionprocess$ is $(\max(1,L_P),L_r)$-LC.
\end{prop}
\begin{proof}
    First, we analyse the reward function.
    Let $x=(s,a_1,\dots,a_\delay)$ and $x'=(s','a_1,\dots,a_\delay')$ in $\augmentedstatespace$ and $a,a'\in\actionspace$ one has,
    \begin{align*}
        \lvert \augmentedexpectedrewardfunction(x,a)-\augmentedexpectedrewardfunction(x',a')\rvert 
        &= \lvert \expectedrewardfunction(s,a_1)-\expectedrewardfunction(s',a_1')\rvert 
        \\
        &\le
        L_r \left( \distance_{\statespace}(s,s') + \distance_{\actionspace}(a_1,a_1')\right)
        \\
        &\le
        L_r \left( \distance_{\augmentedstatespace}(x,x') + \distance_{\actionspace}(a,a')\right),
    \end{align*}
    where we have used \Cref{eq:delayed_reward_deterministic}.
    For the transition function, we note $a_{\delay+1}=a$ and $a_{\delay+1}'=a'$ for convenience. 
    One has,
    \begin{align*}
        \wassersteindistance[1]&(\augmentedtransitionfunction(\cdot\vert x,a_{\delay+1})\Vert \augmentedtransitionfunction(\cdot\vert x',a_{\delay+1}'))
        \\
        &= \sup_{\left\Vert f\right\Vert_L\leq 1} \left\vert \int_{\augmentedstatespace} f(x'')(\augmentedtransitionfunction(x''\vert x,a_{\delay+1})-\augmentedtransitionfunction(x''\vert x',a_{\delay+1}'))\; (dx'') \right\vert
        \\
        &= \sup_{\left\Vert f\right\Vert_L\leq 1} \left\vert \int_{\augmentedstatespace} f(x'')\left[\transitionfunction(s''\vert s,a_{1})\prod_{i=1}^{\delay}\delta_{a_{i+1}}(a_i'')-\transitionfunction(s''\vert s',a_{1}')\prod_{i=1}^{\delay}\delta_{a_{i+1}}(a_i'')\right]\; (dx'') \right\vert
        \\
        &\le \sup_{\left\Vert f\right\Vert_L\leq 1} \left\vert \underbrace{\int_{\augmentedstatespace} f(x'')\prod_{i=1}^{\delay}\delta_{a_{i+1}}(a_i'')\left[\transitionfunction(s''\vert s,a_{1})-\transitionfunction(s''\vert s',a_{1}')\right]\; (dx'')}_{A} \right\vert
        \\
        &\qquad + \sup_{\left\Vert f\right\Vert_L\leq 1} \left\vert \underbrace{\int_{\augmentedstatespace} f(x'')\transitionfunction(s''\vert s',a_{1}') \left[\prod_{i=1}^{\delay}\delta_{a_{i+1}}(a_i'')-\prod_{i=1}^{\delay}\delta_{a_{i+1}}(a_i'')\right]\; (dx'')}_{B} \right\vert,
    \end{align*}     
    where we have added and subtracted the same quantity in the last inequality and used the notation $x''=(s'',a_1'',\dots,a_\delay'')$.
    We consider the first term,
    \begin{align*}
        A &= \int_{\actionspace^\delay} \prod_{i=1}^{\delay}\delta_{a_{i+1}}(a_i'') \int_{\statespace}f(x'')\left[\transitionfunction(s''\vert s,a_{1})-\transitionfunction(s''\vert s',a_{1}')\right]\;(ds'' da_1'' \dots da_{\delay}'')
        \\
        &\le L_P(\distance_{\statespace}(s,s') + \distance_{\actionspace}(a_1,a_1')) \int_{\actionspace^
        {\delay}} \prod_{i=1}^{\delay}\delta_{a_{i+1}}(a_i'')  \;(da_1'' \dots da_{\delay}'')
        \\
        &= L_P(\distance_{\statespace}(s,s') + \distance_{\actionspace}(a_1,a_1')),
    \end{align*}
    where we have used the fact that being 1-LC over $\augmentedstatespace$, $f$ is also 1-LC over $\statespace$.
    For the other term, similarly,
    \begin{align*}
        B &=\int_{\statespace}\transitionfunction(s''\vert s',a_{1}') \int_{\actionspace^\delay}f(x'') \left[\prod_{i=1}^{\delay}\delta_{a_{i+1}}(a_i'')-\prod_{i=1}^{\delay}\delta_{a_{i+1}}(a_i'')\right]\;\;(da_1'' \dots da_{\delay}''ds'')
        \\
        &\leq \sum_{i=2}^{\delay+1}\distance_{\actionspace}(a_{i},a_{i}').
    \end{align*}
    Regrouping the terms, one has,
    \begin{align*}
        \wassersteindistance[1](\augmentedtransitionfunction(\cdot\vert x,a_{\delay+1})\Vert \augmentedtransitionfunction(\cdot\vert x',a_{\delay+1}'))
        &\le L_P(\distance_{\statespace}(s,s') + \distance_{\actionspace}(a_1,a_1')) + \sum_{i=2}^{\delay+1}\distance_{\actionspace}(a_{i},a_{i}')
        \\
        &\le \max(1,L_P) (\distance_{\statespace}(x,x') + \distance_{\actionspace}(a_{\delay+1},a_{\delay+1}'))
    \end{align*}
\end{proof}


\begin{prop}
\label{pp:lip_dida_pol}
    Let $\markovdecisionprocess$ be a $(L_T)$-TLC \gls{mdp}.
    Let $\delayedmarkovdecisionprocess$ be a constantly $\delay$-delayed \gls{dmdp} built upon $\markovdecisionprocess$.
    Assume that a policy $\delayedpolicy$ for $\delayedmarkovdecisionprocess$ satisfies \Cref{eq:belief_pol} for some $L_\pi$-LC policy $\pi$ in $\markovdecisionprocess$.
    Then, $\delayedpolicy$ is $2\delay L_\pi L_T$-LC.
\end{prop}
\begin{proof}
    Let $x=(s,a_1,\dots,a_\delay)$ and $x'=(s','a_1,\dots,a_\delay')$ in $\augmentedstatespace$
    \begin{align}
        \wassersteindistance[1](\delayedpolicy(\cdot\vert x)\Vert \delayedpolicy(\cdot\vert x'))
        &= \sup_{\left\Vert f\right\Vert_L\leq 1} \left\vert \int_{\actionspace} f(a)\pi(a\vert s)(\augmentedbelief(s\vert x)-\augmentedbelief(s\vert x'))\; (da) \right\vert
        \nonumber\\
        &\le L_\pi\sup_{\left\Vert f\right\Vert_L\leq 1} \left\vert \int_{\statespace} f(s)(\augmentedbelief(s\vert x)-\augmentedbelief(s\vert x'))\; (da) \right\vert
        \label{eq:lpi_in_s}
        \\
        &\le 2 L_\pi L_T,
        \label{eq:lpilc}
    \end{align}
    where we have used that $s\mapsto f(a)\pi(a\vert s)$ is $L_pi$-LC in \Cref{eq:lpi_in_s} and \Cref{lem:bound_sigma_tlc} in \Cref{eq:lpilc}.
\end{proof}

\subsection{Bounding \texorpdfstring{$\sigma_b^{\rho}$}{}}
In this sub-section, we provide two ways to bound the term $\textstyle\sigma_b^\rho = \expectedvalue_{\substack{x'\sim \delayeddiscountedstateoccupancydistribution[x][\delayedpolicy]\\ s,s'\sim \augmentedbelief(\cdot\vert x')}}\left[ \distance_{\statespace}(s,s') \right]$ from \Cref{subsec:imitation_upper_bound}, where 
$\rho$ is a distribution on $\augmentedstatespace$. 
For the first one, we assume that $\statespace\subset\realnumbers^n$ is equipped with the Euclidean norm.

\begin{lemma}[Euclidean bound]
\label{lem:bound_sigma_eucl}
    Let $\markovdecisionprocess$ be an \gls{mdp} where $\statespace\subset\realnumbers^n$ is equipped with the Euclidean norm. Then, one has
    \begin{align*}
        \sigma_b^\rho \leq \sqrt{2}\expectedvalue_{x'\sim \delayeddiscountedstateoccupancydistribution[\rho][\delayedpolicy](\cdot)}\left[\sqrt{ \variance_{s\sim \augmentedbelief(\cdot|x')}(s|x')}\right],
    \end{align*}
    where $\rho$ is a distribution on $\augmentedstatespace$.
\end{lemma}
\begin{proof}
    The proof stems from the following steps,
    \begin{align}
        \sigma_b^\rho 
        &= \expectedvalue_{\substack{x'\sim \delayeddiscountedstateoccupancydistribution[\rho][\delayedpolicy]\\ s,s'\sim \augmentedbelief(\cdot\vert x')}}\left[ \distance_{\statespace}(s,s') \right]\nonumber
        \\
        &= \expectedvalue_{\substack{x'\sim \delayeddiscountedstateoccupancydistribution[\rho][\delayedpolicy]\\ s,s'\sim \augmentedbelief(\cdot\vert x')}}\left[ \sqrt{(s'-s)^2} \right]\label{eq:def_euclid}
        \\
        &= \expectedvalue_{x'\sim \delayeddiscountedstateoccupancydistribution[\rho][\delayedpolicy]}\sqrt{\expectedvalue_{s,s'\sim \augmentedbelief(\cdot\vert x')} \left[ (s'-s)^2 \right]}\label{eq:jensen_var}
        \\
        &= \expectedvalue_{x'\sim \delayeddiscountedstateoccupancydistribution[\rho][\delayedpolicy]}\left[\sqrt{ \variance_{s\sim \augmentedbelief(\cdot|x')}(s|x')}\right].\label{eq:idd_var_s},
    \end{align}
    where in \Cref{eq:def_euclid} the Euclidean norm is explicited, in \Cref{eq:jensen_var} Jensen's inequality is applied and 
    \Cref{eq:idd_var_s} follows from the fact that, if $s'$ and $s$ are i.i.d., one has 
    $\textstyle\expectedvalue_{s,s'\sim \augmentedbelief(\cdot\vert x')} \left[ (s'-s)^2 \right] =
    2\variance_{s\sim \augmentedbelief(\cdot|x')}[s]$.
\end{proof}


Let us now provide the second result, which assumes time-Lipschitzness of the \gls{mdp} (see \Cref{def:time_lip} and \Cref{eq:time_lip_non_int}).
Before providing the result, we prove the following intermediate result.

\begin{prop}
\label{pp:tlc_delay}
    Let $\markovdecisionprocess$ be an $L_T$-TLC \gls{mdp} \footnote{Note that, since non-integer delays are considered, we extended the TLC assumption in \Cref{eq:time_lip_non_int}}. Let $x=(s_1,a_1,\dots,a_\delay)\in\statespace\times\actionspace^\delay$ be an augmented state for a delay $\delay\in\realnumbers[\ge0]$.
    Then, one has:
    \begin{align*}
        \wassersteindistance\left(\augmentedbelief(\cdot\vert x)\Vert \delta_{s_1} \right)\leq \delay L_T
    \end{align*}
\end{prop}
\begin{proof}
    We prove the result first for integer delays then for non-integer delays.
    
    \noindent\textbf{Integer delay.}\indent Let $\delay\in\naturalnumbers$.
    The proof is made by induction. 
    For the initialization, when $\delay=0$, the result is clearly valid because the current state is known to the agent. 
    Note that $\delay=1$ holds by the $L_T$-TLC assumption.
    For the recurrence, assume now that the statement holds for some $\delay \in\naturalnumbers$. Let $x=(s_1,a_1,\dots.a_{\delay+1})\in\statespace\times\actionspace^{\delay+1}$. Then,
    \begin{align}
        &\wassersteindistance\left(\augmentedbelief[\delay+1](\cdot\vert x)\Vert \delta_{s_1} \right) 
        \\
        &= \sup_{\left\Vert f\right\Vert_L\leq 1}\left\vert\int_{\statespace} f(s')\left( \augmentedbelief[\delay+1](s' \vert x) -\delta_{s_1}(s')\right)ds' \right\vert\nonumber
        \\
        &= \sup_{\left\Vert f\right\Vert_L\leq 1}\left\vert\int_{\statespace} p(s_2\vert s_1,a_1) \int_{\statespace} f(s')\left( \augmentedbelief(s' \vert s_2,a_2,\cdots,a_\delay) -\delta_{s_1}(s')\right)ds' \right\vert
        \label{eq:cond_s2}
        \\
        &= \sup_{\left\Vert f\right\Vert_L\leq 1}\left\vert\int_{\statespace} p(s_2\vert s_1,a_1) \int_{\statespace} f(s')\left( \augmentedbelief(s' \vert s_2,a_2,\cdots,a_\delay) -\delta_{s_2}(s')\right)ds' \right.\nonumber
        \\
        &\qquad + \left.\int_{\statespace} p(s_2\vert s_1,a_1) \int_{\statespace} f(s')\left( \delta_{s_2}(s) -\delta_{s_1}(s')\right)ds' \right\vert
        \label{eq:add_remove_delta}
        \\
        &\leq \underbrace{\sup_{\left\Vert f\right\Vert_L\leq 1}\left\vert\int_{\statespace} p(s_2\vert s_1,a_1) \int_{\statespace} f(s')\left( \augmentedbelief(s' \vert s_2,a_2,\cdots,a_\delay) -\delta_{s_2}(s')\right)ds'  \right\vert}_A\nonumber
        \\
        &\qquad + \underbrace{\wassersteindistance\left(P(\vert s_1,a_1)\Vert \delta_{s_1}\right)}_B
        \nonumber,
    \end{align}
    where \eqref{eq:cond_s2} is obtained by conditionning on the next observed state $s_2$ and \Cref{eq:add_remove_delta} is obtained by adding and subtracting the same quantity $\delta_{s_2}$.
    
    The term $A$ can be bounded using the statement at $\delay$ while $B$ falls directly under the TLC assumption. 
    We therefore have 
    \begin{align*}
        \wassersteindistance\left(\augmentedbelief[\delay+1](\cdot\vert x)\Vert \delta_{s_1} \right) \leq (\delay+1) L_T.
    \end{align*}
    By induction, we proved the result for any $\delay\in\naturalnumbers$.
    
    \noindent\textbf{Non-integer delay.}\indent To show this result in the more general case of $\delay\in\realnumbers[\ge0]$, one can divide the delay in its fractional ($\fractionalpart{\delay}$) and integer part ($\integerpart{\delay}$). Applying a similar reasoning as before, one would get similar terms $A$ and $B$ where $A$ involves the integer part of the delay for which the statement holds and $B$ involves the fractional part of $\delay$ for which the statement holds under \Cref{eq:time_lip_non_int}. 
\end{proof}



\begin{lemma}[Time-Lipschitz bound]
\label{lem:bound_sigma_tlc}
    Let $\markovdecisionprocess$ be an $L_T$-TLC \gls{mdp}. Consider the $\delay$-delayed \gls{dmdp} $\delayedmarkovdecisionprocess$ obtained from $\markovdecisionprocess$ for $\delay\in\realnumbers[\ge0]$. Then,
    \begin{align*}
        \sigma_b^\rho  \leq 2\delay L_T,
    \end{align*}
    where $\rho$ is a distribution on $\augmentedstatespace$.
\end{lemma}
\begin{proof}
    For an augmented state $x$, let $s_x$ be the last observed state it contains.
    Then we can write the following,
    \begin{align}
        \sigma_b^\rho
        &= \expectedvalue_{\substack{x\sim \delayeddiscountedstateoccupancydistribution[\rho][\delayedpolicy]\\ s,s'\sim \augmentedbelief(\cdot\vert x)}}\left[ \distance_{\statespace}(s,s') \right]\nonumber
        \\
        &\leq \expectedvalue_{\substack{x\sim \delayeddiscountedstateoccupancydistribution[\rho][\delayedpolicy]\\ s\sim \augmentedbelief(\cdot\vert x')}}\left[ \distance_{\statespace}(s,s_x) \right] + \expectedvalue_{\substack{x\sim \delayeddiscountedstateoccupancydistribution[\rho][\delayedpolicy]\\ s'\sim \augmentedbelief(\cdot\vert x)}}\left[ \distance_{\statespace}(s_x,s') \right]\label{eq:triang_ineq_tlc}
        \\
        &= 2\expectedvalue_{x\sim \delayeddiscountedstateoccupancydistribution[\rho][\delayedpolicy]}\left[\int_{\statespace} d_{\statespace}(s,s_x) \augmentedbelief(s\vert x)~ds\right] \nonumber
        \\
        &= 2\expectedvalue_{x\sim \delayeddiscountedstateoccupancydistribution[\rho][\delayedpolicy]}\left[\int_{\statespace} \distance_{\statespace}(s,s_x) \left(\augmentedbelief(s\vert x) - \delta_{s_x}(s) \right)~ds\right]\label{eq:null_term}
        \\
        &\leq 2 \expectedvalue_{x\sim \delayeddiscountedstateoccupancydistribution[\rho][\delayedpolicy]}\left[\wassersteindistance\left(\augmentedbelief(\cdot\vert x)\Vert \delta_{s_x} \right)\right]\label{eq:wass_recognise},
    \end{align}
    where \Cref{eq:triang_ineq_tlc} holds by triangular inequality; in \Cref{eq:null_term}, the term  $\textstyle\int_{\statespace}d_{\statespace}(s,s_x)\delta_{s_x}(s)ds=0$ is added; in \Cref{eq:wass_recognise} the definition of the Wasserstein distance is used.
    To conclude, \Cref{pp:tlc_delay} is applied within the expectation.
\end{proof}